
\documentclass[12pt]{article}

\usepackage[margin=1in]{geometry}
\usepackage{longtable}
\usepackage{hyperref}
\usepackage{enumitem}

\title{Scientific Forgery Image Detection\\SRS Version 2}
\author{Janelle Rivera \and Eric Marroquin \and Alexis Flores \and Alex Lam}
\date{May 8, 2026}

\begin{document}

\maketitle

\tableofcontents
\newpage

\section{Introduction}

\subsection{Purpose}

This document defines the requirements for the Scientific Forgery Image Detection
System. The system is designed to detect copy-move forgery in biomedical and
scientific images, where regions of an image are duplicated to manipulate research
results.

The system will allow users to upload scientific images through a web interface. These
images will be processed by a backend system designed to support future detection
functionality.

\subsection{Intended Audience}

This document is intended for:

\begin{itemize}
    \item Software developers working on the system
    \item Course instructors evaluating the project
    \item Students involved in system design
    \item Future maintainers of the software
\end{itemize}

\subsection{Overall Description}

The system is a web-based application consisting of:

\begin{itemize}
    \item A frontend user interface for uploading images
    \item A backend API for receiving and processing images
    \item A planned detection system for analyzing scientific images
\end{itemize}

At this stage, the system is in early development and primarily focuses on:

\begin{itemize}
    \item Project setup
    \item Image upload functionality
    \item Backend communication
    \item Basic workflow structure
\end{itemize}

No advanced image detection model is implemented yet.

\section{External Interface Requirements}

\subsection{User Interface Requirements}

The system will provide a simple web interface with the following components:

\begin{itemize}
    \item Homepage displaying project information
    \item Image upload page
    \item Upload button for selecting image files
    \item Submit button to send image to backend
    \item Basic result display area
\end{itemize}

The interface must be simple, accessible, and user-friendly.

\subsection{Software Interface Requirements}

The system will interact with the following software components:

\textbf{Frontend:}

\begin{itemize}
    \item React.js for UI development
    \item HTML/CSS for layout and styling
\end{itemize}

\textbf{Backend:}

\begin{itemize}
    \item Python as the programming language
    \item FastAPI for API development
    \item Uvicorn as the server runtime
\end{itemize}

\textbf{Communication:}

\begin{itemize}
    \item HTTP requests between frontend and backend
    \item JSON responses from backend
\end{itemize}

\subsection{Hardware Interface Requirements}

No specialized hardware is required. The system will run on standard computing
devices capable of running a web browser and local development servers.

\subsection{Communication Requirements}

\begin{itemize}
    \item REST API communication between frontend and backend
    \item Image files transmitted via HTTP requests
    \item Responses returned in JSON format
\end{itemize}

\section{Legal and Ethical Considerations}

\subsection{Data Integrity}

The system processes scientific images, which may contain sensitive or
research-related content. All uploaded images should be handled securely and not
altered except for analysis purposes.

\subsection{Privacy Considerations}

\begin{itemize}
    \item Uploaded images are used only for detection purposes
    \item No user personal data is required
    \item Data should not be shared externally without permission
\end{itemize}

\subsection{Ethical Concerns}

This system is designed to support scientific integrity by detecting possible image
manipulation. However:

\begin{itemize}
    \item Results should not be considered absolute proof of fraud
    \item The system should be used as a supporting tool
    \item Misinterpretation of results must be avoided
\end{itemize}

\section{Glossary}

\begin{longtable}{|p{4cm}|p{9cm}|}
\hline
\textbf{Term} & \textbf{Definition} \\
\hline
API & Interface used for communication between software components \\
\hline
Backend & Server-side system that processes requests \\
\hline
Frontend & User interface of the application \\
\hline
Copy-Move Forgery & Image manipulation technique where parts of an image are duplicated \\
\hline
FastAPI & Python framework for building APIs \\
\hline
React.js & JavaScript library for building user interfaces \\
\hline
HTTP & Protocol used for web communication \\
\hline
JSON & Lightweight data format used for API responses \\
\hline
\end{longtable}

\section{Versions Table}

\begin{longtable}{|p{1.5cm}|p{2.5cm}|p{5.5cm}|p{4cm}|}
\hline
\textbf{Version} & \textbf{Date} & \textbf{Description} & \textbf{Authors} \\
\hline
1.0 & May 8, 2026 & Initial SRS created for Snapshot 1 & Janelle Rivera, Eric Marroquin, Alexis Flores, Alex Lam \\
\hline
1.1 & May 10, 2025 & Updated to include frontend implementation progress, UI development, and revised system interface requirements & Janelle Rivera, Eric Marroquin, Alexis Flores, Alex Lam \\
\hline
\end{longtable}

\section{References}

\begin{itemize}
    \item FastAPI Documentation: \url{https://fastapi.tiangolo.com/}
    \item React Documentation: \url{https://react.dev/}
    \item OWASP API Security Guidelines: \url{https://owasp.org/www-project-api-security/}
\end{itemize}

\end{document}
